% GENERATED FILE. Edit canonical lessons, then run npm run generate:books.
% canonical-source-hash: fnv1a64:f692944483ea3efc
% canonical-lessons: ZH-C186-jiandao, ZH-C186-lajitong, ZH-C186-suliaodai, ZH-C186-pingzi, ZH-C186-chouti

\chapter{Odds and Ends}
\label{ch:zh-odds-and-ends}
\hlchaptermodality{186}

\begin{chapteropening}
\textbf{By the end of this chapter:} \emph{I can name scissors, a rubbish bin, a plastic bag, a bottle and a drawer.}

Complete the last lesson of chapter 186: I can name scissors, a rubbish bin, a plastic bag, a bottle and a drawer.
\end{chapteropening}

\section[jiǎndāo]{jiǎndāo — scissors}

\label{lesson:ZH-C186-jiandao}



\begin{quote}
Before the new one: say the Mandarin for morning and evening, then the Mandarin for more often than not.
\end{quote}

\subsection*{You'll want to know: jiǎndāo}
\emph{jiǎndāo} — \textquotedblleft{}scissors\textquotedblright{}.

\begin{quote}
\emph{jiǎndāo zài nǎr} — Where are the scissors?
\end{quote}

Its characters come later; for now it is a word for the ear.

\begin{cousinweb}
\emph{jiǎndāo} is built from two parts inherited from Middle Chinese (Guangyun, in Baxter's transcription): \emph{tsjenX}, glossed 'to cut off' and 'to make even', and \emph{taw}, the word for a knife: a cutting knife.
\end{cousinweb}

\begin{grammarlens}[title={What you've built}]
One more word for everyday things.
\end{grammarlens}

\subsection*{Your turn}
\begin{itemize}
\raggedright
  \item \emph{Say it:} \emph{jiǎndāo}
  \item \emph{Say it:} \emph{jiǎndāo}, once more
  \item \emph{Say it:} \emph{jiǎndāo zài nǎr} — Where are the scissors?
  \item \emph{Recall:} say the Mandarin for morning and evening, then the Mandarin for more often than not, then say \emph{jiǎndāo} again
\end{itemize}

\begin{tcolorbox}[breakable,colback=teal!4,colframe=teal!35!black,title={Before you move on}]
What does jiǎndāo mean? (\textquotedblleft{}Scissors\textquotedblright{}.) Say it once more.
\end{tcolorbox}
\section[lājītǒng]{lājītǒng — rubbish bin}

\label{lesson:ZH-C186-lajitong}



\begin{quote}
Before the new one: say the Mandarin for more often than not, then the Mandarin for scissors.
\end{quote}

\subsection*{You'll want to know: lājītǒng}
\emph{lājītǒng} — \textquotedblleft{}rubbish bin\textquotedblright{}.

\begin{quote}
\emph{lājītǒng mǎn le} — The bin is full.
\end{quote}

Its characters come later; for now it is a word for the ear.

\begin{cousinweb}
\emph{lājītǒng} is built from \emph{lājī} 'rubbish' and \emph{tǒng} 'tub'. The second is inherited from Middle Chinese \emph{thuwngX} (Guangyun, in Baxter's transcription), glossed 'a wooden tub'; for \emph{lājī} no dictionary we use records an older form, so the history of that part is not settled.
\end{cousinweb}

\begin{grammarlens}[title={What you've built}]
One more word for everyday things.
\end{grammarlens}

\subsection*{Your turn}
\begin{itemize}
\raggedright
  \item \emph{Say it:} \emph{lājītǒng}
  \item \emph{Say it:} \emph{lājītǒng}, once more
  \item \emph{Say it:} \emph{lājītǒng mǎn le} — The bin is full.
  \item \emph{Recall:} say the Mandarin for more often than not, then the Mandarin for scissors, then say \emph{lājītǒng} again
\end{itemize}

\begin{tcolorbox}[breakable,colback=teal!4,colframe=teal!35!black,title={Before you move on}]
What does lājītǒng mean? (\textquotedblleft{}Rubbish bin\textquotedblright{}.) Say it once more.
\end{tcolorbox}
\section[sùliàodài]{sùliàodài — plastic bag}

\label{lesson:ZH-C186-suliaodai}



\begin{quote}
Before the new one: say the Mandarin for scissors, then the Mandarin for a rubbish bin.
\end{quote}

\subsection*{You'll want to know: sùliàodài}
\emph{sùliàodài} — \textquotedblleft{}plastic bag\textquotedblright{}.

\begin{quote}
\emph{wǒ bú yào sùliàodài} — I don't need a plastic bag.
\end{quote}

Its characters come later; for now it is a word for the ear.

\begin{cousinweb}
\emph{sùliàodài} is built from \emph{sùliào} 'plastic' and \emph{dài} 'bag'. Of its parts, \emph{liào} 'material' and \emph{dài} are inherited from Middle Chinese, \emph{lew} and \emph{dojH} (Guangyun, in Baxter's transcription); the first is listed there as \emph{suH}, glossed 'a modelled figure': a material that can be modelled.
\end{cousinweb}

\begin{grammarlens}[title={What you've built}]
One more word for everyday things.
\end{grammarlens}

\subsection*{Your turn}
\begin{itemize}
\raggedright
  \item \emph{Say it:} \emph{sùliàodài}
  \item \emph{Say it:} \emph{sùliàodài}, once more
  \item \emph{Say it:} \emph{wǒ bú yào sùliàodài} — I don't need a plastic bag.
  \item \emph{Recall:} say the Mandarin for scissors, then the Mandarin for a rubbish bin, then say \emph{sùliàodài} again
\end{itemize}

\begin{tcolorbox}[breakable,colback=teal!4,colframe=teal!35!black,title={Before you move on}]
What does sùliàodài mean? (\textquotedblleft{}Plastic bag\textquotedblright{}.) Say it once more.
\end{tcolorbox}
\section[píngzi]{píngzi — bottle}

\label{lesson:ZH-C186-pingzi}



\begin{quote}
Before the new one: say the Mandarin for a rubbish bin, then the Mandarin for a plastic bag.
\end{quote}

\subsection*{You'll want to know: píngzi}
\emph{píngzi} — \textquotedblleft{}bottle\textquotedblright{}.

\begin{quote}
\emph{zhège píngzi lǐ yǒu shuǐ} — There's water in this bottle.
\end{quote}

Its characters come later; for now it is a word for the ear.

\begin{cousinweb}
\emph{píngzi} is built from \emph{píng} and the noun ending \zh{子}, both inherited from Middle Chinese, \emph{beng} and \emph{tsiX} (Guangyun, in Baxter's transcription), which glosses the first as a vessel for drawing water. In the word the second loses its tone.
\end{cousinweb}

\begin{grammarlens}[title={What you've built}]
One more word for everyday things.
\end{grammarlens}

\subsection*{Your turn}
\begin{itemize}
\raggedright
  \item \emph{Say it:} \emph{píngzi}
  \item \emph{Say it:} \emph{píngzi}, once more
  \item \emph{Say it:} \emph{zhège píngzi lǐ yǒu shuǐ} — There's water in this bottle.
  \item \emph{Recall:} say the Mandarin for a rubbish bin, then the Mandarin for a plastic bag, then say \emph{píngzi} again
\end{itemize}

\begin{tcolorbox}[breakable,colback=teal!4,colframe=teal!35!black,title={Before you move on}]
What does píngzi mean? (\textquotedblleft{}Bottle\textquotedblright{}.) Say it once more.
\end{tcolorbox}
\section[chōuti]{chōuti — drawer}

\label{lesson:ZH-C186-chouti}



\begin{quote}
Before the new one: say the Mandarin for a plastic bag, then the Mandarin for a bottle.
\end{quote}

\subsection*{You'll want to know: chōuti}
\emph{chōuti} — \textquotedblleft{}drawer\textquotedblright{}.

\begin{quote}
\emph{jiǎndāo zài chōuti lǐ} — The scissors are in the drawer.
\end{quote}

Its characters come later; for now it is a word for the ear.

\begin{cousinweb}
\emph{chōuti} is built from \emph{chōu} 'to pull out' and \emph{ti} 'a tray'. The first is inherited from Middle Chinese \emph{trhjuw} (Guangyun, in Baxter's transcription), glossed 'to pull out'; for the second no dictionary we use records an older form, so its history is not settled. In the word it loses its tone.
\end{cousinweb}

\begin{grammarlens}[title={What you've built}]
One more word for everyday things.
\end{grammarlens}

\subsection*{Your turn}
\begin{itemize}
\raggedright
  \item \emph{Say it:} \emph{chōuti}
  \item \emph{Say it:} \emph{chōuti}, once more
  \item \emph{Say it:} \emph{jiǎndāo zài chōuti lǐ} — The scissors are in the drawer.
  \item \emph{Recall:} say the Mandarin for a plastic bag, then the Mandarin for a bottle, then say \emph{chōuti} again
\end{itemize}

\begin{tcolorbox}[breakable,colback=teal!4,colframe=teal!35!black,title={Before you move on}]
What does chōuti mean? (\textquotedblleft{}Drawer\textquotedblright{}.) Say it once more.
\end{tcolorbox}
